Robots that learn from human preference comparisons ask people which short
behaviors look better in order to infer a reward function. Much of the
uncertainty about that reward never changes what the robot should actually do,
so limited human patience is spent resolving differences with no behavioral
consequence.
We select trajectory comparisons by expected reduction in policy regret under a
population prior over previous users, rather than by parameter information gain,
and evaluate on a tabular gridworld with exact optimal policies and exact regret
using simulated mixture-population users.
A population particle prior substantially reduces final policy regret relative to
a uniform particle prior under the same VOI acquisition
(C-mech-2; E-prior), and a misspecified prior hurts
(C-mech-3; E-prior-shift).
Relative to BALD, one-step VOI is not yet a supported win on final regret
(C-main-2; E3-power / E3-power-120 / E-decouple-power);
VOI remains only fragilely better than random (C-main-1), with a large remaining
gap to a true-regret oracle (C-main-3).
Changing the acting rule (map / sample / softminimax) does not close that gap
(C-mech-4; E-act).
These results support using a population prior for decision-relevant elicitation
while clarifying that beating parameter-uncertainty acquisition remains open.
% Cite Regan \& Boutilier 2009 and Model-Free Preference Elicitation (IJCAI 2024) in the intro.
